\documentclass{tufte-handout}

%\geometry{showframe}% for debugging purposes -- displays the margins

\usepackage[spanish, es-tabla]{babel}
\usepackage{amsmath}

% Set up the images/graphics package
\usepackage{graphicx}
\setkeys{Gin}{width=\linewidth,totalheight=\textheight,keepaspectratio}
\graphicspath{{graphics/}}

\title[Química Física II: Ejercicios Tema 2]{
Química Física II: Ejercicios Tema 2}
%\author{David De Sancho}
\date{}  % if the \date{} command is left out, the current date will be used

% The following package makes prettier tables.  We're all about the bling!
\usepackage{booktabs}

% The units package provides nice, non-stacked fractions and better spacing
% for units.
\usepackage{units}

% The fancyvrb package lets us customize the formatting of verbatim
% environments.  We use a slightly smaller font.
\usepackage{fancyvrb}
\fvset{fontsize=\normalsize}

% Small sections of multiple columns
\usepackage{multicol}

% Provides paragraphs of dummy text
\usepackage{lipsum}

% These commands are used to pretty-print LaTeX commands
\newcommand{\doccmd}[1]{\texttt{\textbackslash#1}}% command name -- adds backslash automatically
\newcommand{\docopt}[1]{\ensuremath{\langle}\textrm{\textit{#1}}\ensuremath{\rangle}}% optional command argument
\newcommand{\docarg}[1]{\textrm{\textit{#1}}}% (required) command argument
\newenvironment{docspec}{\begin{quote}\noindent}{\end{quote}}% command specification environment
\newcommand{\docenv}[1]{\textsf{#1}}% environment name
\newcommand{\docpkg}[1]{\texttt{#1}}% package name
\newcommand{\doccls}[1]{\texttt{#1}}% document class name
\newcommand{\docclsopt}[1]{\texttt{#1}}% document class option name

\begin{document}
\maketitle
\large
\begin{enumerate}
    \item Suponiendo que $\psi_1$ y
    $\psi_2$ son dos funciones reales que 
    están normalizadas y son ortogonales, 
    normaliza las siguientes funciones:
    (a) $\psi_1+\psi_2$, (b) 
    $\psi_1-\psi_2$.
    
    \item ¿Cuál de las siguientes funciones
    es función propia del operador
    $\hat{D}=\frac{d}{dx}$? (a) $kx^2$
    (b) $\sin x$ (c) $\exp (kx)$ 
    (d) $\exp(kx^2)$ (e) $\exp(ikx)$.
    
    \item Los operadores para la posición y
    el momento lineal vienen dados por 
    \begin{equation}
        \begin{array}{cc}
             \hat{x}& =x  \\
            \hat{p}_x& =\frac{\hbar}{i}\big(\frac{\partial}{\partial x}\big)
        \end{array}
    \end{equation}
    (a) ¿Son estos operadores hermíticos?
    (b) ¿Conmutan estos operadores?
    (c) ¿Son los operadores citados lineales
    o no lineales?
    (d) ¿Es la función $\psi_1=A\sin (n\pi
    x/a)$, donde $A$, $n$ y $a$ son
    constantes, una función propia de ambos
    operadores?
   
   \item Normaliza la función del ejercicio
   3d en el dominio $0\leq x\leq a$.
   
   \item Determina las condiciones necesarias 
   para que la función $\psi(x)=x\exp(-\alpha
   x^2)$ sea función propia del operador 
   $\hat{\Omega}=\frac{d^2}{dx^2} + bx^2$.
   Normalizar la función y calcular el valor 
   propio de $x$ para un
   sistema descrito por dicha función. Para resolver
   este ejercicio es necesaria la integral
   \begin{equation}
       \int_0^{\infty}x^{2n}e^{-ax^2}dx = 
       \frac{1\cdot 3\cdot 5... (2n-1)}{2^{n+1}a^n}\sqrt{\frac{\pi}{a}}
   \end{equation}

   \item Demuestra que la combinación lineal de
   operadores $\hat{A}+\mathrm{i}\hat{B}$ no es hermítica
   si $\hat{A}$ y $\hat{B}$ son hermíticos.

   \item Determina el conmutador de los operadores
   $\hat{a}$ y $\hat{a}^\dagger$, donde $\hat{a}=2^{-1/2}(\hat{x}+\mathrm{i}\hat{p})$
   y $\hat{a}^\dagger=2^{-1/2}(\hat{x}-\mathrm{i}\hat{p})$.
\end{enumerate}

\newpage
\subsection{\textbf{Soluciones:}}
\begin{enumerate}
    \item (a) $\psi=(\frac{1}{\sqrt{2}})(\psi_1+\psi_2)$. (b) $\psi=(\frac{1}{\sqrt{2}})(\psi_1-\psi_2)$
    \item (a) No. (d) No. (c) Sí. (d) No. (e) Sí.
    \item (a) Tanto $\hat{x}$ como $\hat{p}_x$ son hermíticos. (b) No conmutan. (c) Ambos son lineales. (d) La función $\psi$ no es función propia ni de $\hat{x}$ ni de $\hat{p}_x$.
    \item La constante de normalización es $A=\sqrt{2/a}$.
    \item La función $\psi$ es función propia si $b=-4\alpha^2$. La constante de normalización es
    $N=(\frac{2}{\sqrt{\pi/(2a)^3}})^{1/2}$. El valor
    promedio de la posición es $\langle x \rangle=0$ (la función que debemos integrar es una función impar).
    \item --
    \item $[\hat{a},\hat{a}^\dagger]=\hbar$
\end{enumerate}
\end{document}